\cotefile{156-4}
\scanpage{58}\mspage{57}
\diagram{\begin{tikzcd}[ampersand replacement=\&,column sep=large]
\underset{\substack{\text{ss-fig de}\\ F}}{\mathcal F_{\leq F}}
 \times
\underset{\substack{\text{ss fig}\\ \text{de }G}}{\mathcal F_{\leq G}}
 \arrow[r,"\sim"] \&
\underset{\substack{\text{ss. fig}\\ \text{de }F\cup G=H}}{\mathcal F_{\leq F\cup G}}
\end{tikzcd}}

(On récupère $F',G'$ \unsure{:|à} partir de $H'=F'\cup G'$ comme $H\cap F$ et $H\cap G$.) \unsure{Item|Idem} pour \unsure{une} \deleted{\unsure{familles|famille} \unsure{mutuellement} disjointes d'une figure (\unsure{y}} \inserted{des figures deux à deux disjointes \unsure{fais.t une|ds une}} configuration (famille \unsure{pas|par} \unsure{néc.t|néct} finie).

\uline{Démonstration} \unsure{Particulièrt|Particulièrmt} claire, quand on interprète $\mathcal F$ en termes de l'ens $\unsure{\mathcal L|L}\unsure{\subseteq|\subset}\mathcal P(M)$ satisfaisant à a) b) c). L'ens ordonné $M_{\leq F\vee G}$ s'identifie à la réunion disjointe dans l'ens ordonné somme de $M_{\leq F}$ et $M_{\leq G}$, donc ses parties fermées \deleted{\unsure{sont}} \unsure{s'en|sont} \unsure{définies|définissent|déduisent} par les couples de parties fermées de $M_{\leq F}$ et $M_{\leq G}$ \ldots

\uline{Remarque}
% slanted
\margin{à joindre à la remarque de la page 55}
En termes de l'ens \uline{ordonné} $M$, la donnée de $\mathcal L$ équivaut \unsure{à|:} un ens de parties \uline{fermées} de $M$, satisfaisant
\[
\begin{aligned}
\text{a) }&\deleted{\unsure{\mathcal L}}\forall\,\deleted{\unsure{=}}A\in\mathcal L,\quad
 \deleted{\text{\unsure{toute}}}\inserted{\text{et}}\ \text{partie fermée }A'\text{ de }A,\ \text{on a }A'\in\mathcal L\\
&\deleted{\text{b) }\varnothing\in\mathcal L\quad
 (\Longleftrightarrow\mathcal L\ne\varnothing,\ \text{\unsure{conséquence} de a), et}}\\
\text{b) }&\bigcup_{A\in\mathcal L} A=M\quad\deleted{\unsure{\forall x}}\quad
 \left[\,\underset{\substack{\text{à cause}\\\text{de a)}}}{\Longleftrightarrow}\
 \forall x\in M,\ M_{\leq x}\in\mathcal L\,\right]\\
\text{c) }&\varnothing\in\mathcal L\quad
 \bigl(\underset{\deleted{\unsure{\text{a) b)}}}}{\Longleftrightarrow}
 \mathcal L\ne\varnothing,\ \text{et \unsure{automatiqt} satisfaite}\ \deleted{\unsure{\mathcal L}}\\
 &\hspace{7em}\text{(à cause de b)) si }M\ne\varnothing\bigr)
\end{aligned}
\]

\uline{NB} Soit $\mathcal P_f(M)$ l'ens des parties fermées de $M$, $\mathcal P_{f'}(\mathcal F)$ de $\mathcal F$ qui sont stables par \unsure{Sup|sup}. Si $\mathcal F'$ et $M'$ se correspondent \inserted{\unsure{resp|ssys}}, 
\diagram{\fbox{\begin{tikzcd}[ampersand replacement=\&,column sep=large]
\mathcal P_f(\unsure{\mathcal F|M\mathcal F}) \arrow[r,shift left=.6ex,"\varphi"] \& \mathcal P_f(M) \arrow[l,shift left=.6ex,"\psi"]
\end{tikzcd}}}
$\varphi(\mathcal F')=\mathcal F'\cap M$, $\psi(M')=\{F\in\mathcal F\mid\overline F\unsure{\subset|\not\subset}M'\}$ alors $\varphi,\psi$ bij. inverses l'une de l'autre
\scanpage{59}\mspage{58}
% slanted
\margin{alors $\mathcal F'$ satisfait aussi $C1$, $C2$, et $M'$ \deleted{\unsure{s'identifie}} \inserted{\unsure{désignant|désigne|désignant par}} ses objets irréductibles, la relation d'ordre \unsure{(de|de} $M'\unsure{)|{}}$ induite par $\mathcal F'$ étant celle \deleted{\unsure{d'}}\unsure{induit|induite} par $M$. Si $\mathcal F$ satisfait $C3$, alors $\mathcal F'$ aussi.}

\deleted{\mbox{\uline{C 3}} Soient $F,G,H$} \deleted{trois} \deleted{figures entre elles compatibles.} \deleted{\mbox{\inserted{i.e. deux à deux \unsure{majorés.|majorées.}}}} \deleted{Alors $\{F,G,H\}$ est une configuration i.e.} \deleted{$F,G,H$} \deleted{$\exists$ figure $K$ telle que $F,G,H$ \unsure{sont} des sous figures de $K$ (donc existe borne supérieure de}

\uline{C 3} Soient $F,G\in\mathcal F$, telles que
\[
\forall\,x\in\overline F,\ \unsure{\overleftarrow{y}|\bar y|y}\in\overline G,\quad x\text{ et }y\text{ compatibles}
\]
Alors $F$ et $G$ sont compatibles

% slanted
\margin{\uline{Cor} Pour que $F,G$ disjointes, il f et s. que $\forall x\in\overline F,y\in\overline G$, $x$ et $y$ le soient}

\uline{NB} L'inverse évident : Si $F$ et $G$ sont compatibles, alors pour toutes sous figures $F'$ de $F$ (\unsure{p. ex.|p. e.} une \unsure{ss-fig.|ss fig.} \unsure{él.|ir.}) et $G'$ de $G$ (idem) $F'$ et $G'$ le sont.

\uline{Cor} Soient $F$ \unsure{et} $G,H$ \deleted{\unsure{deux|trois}} figures, \deleted{\unsure{Si}} \unsure{avec} $F,G$ compatibles. Pour que $H$ soit compatible \unsure{:|à} $F\cup G$, il f et s qu'il le soit \unsure{:|à} $F$ et \unsure{à|:} $G$. \deleted{\unsure{Deux}}

Plus gén, soit $(F_i)_{i\in I}$ une configuration, $F$ la \deleted{\unsure{leur|borne}} \inserted{figure} réunion. Pour que $H$ soit compatible \unsure{à|:} $F$, il f et s que $H$ soit compatible avec $F_i$.
\scanpage{60}\mspage{59}
Donc, en \deleted{\unsure{de|te}}termes d'une structure portée par un ens ordonné $M$ (cf remarque page 58, on trouve ceci. On a sur \deleted{\uline{\unsure{Définition}}} $M$ une relation $R$ (``compatibilité) qui est \unsure{réflexive|réflexixe} symétrique (donc équivaut à la donnée de $\unsure{\mathscr S|\mathcal S}\subset\mathcal P_2(M)$), \deleted{\uline{\unsure{impliquant}}}\inserted{et telle que} $R(x,y)$ implique $R(x',y')$ pour $x'\leq x$, $y'\leq y$. \deleted{\unsure{comp stable} par $\leq$.} Cela permet, à peu de choses près, de reconstituer $(\mathcal F,\leq)$, en \unsure{tant qu'}ens de parties de $M$. Pour l'expliciter,

Quelques \unsure{terminologies|terminologie}

``\uline{Composante}'' \uline{irréductible} \inserted{(ou strate maximale)} d'une figure $F$ : \deleted{\unsure{o|a}}strate maximale de $F$ ($\exists$ tjrs si l'ens $\overline F$ des strates de $F$ est fini \inserted{(on dit alors que $F$ est une \uline{figure finie})} et $F\ne\varnothing_{\mathcal F}$).

Figure \deleted{de type}\inserted{de \unsure{type fini|type}} fini : qui est borne sup. de \unsure{l'ens|la réu} de ses ``composantes'' irréductibles (\unsure{exemple|exemples} : toute figure finie est de t.f.), ou encore, toute figure \inserted{$F$} qui est \deleted{de la forme $\operatorname{Sup}\unsure{X|\mathcal X|A}$}\inserted{borne supr. d'un \unsure{nb} \uline{fini} de} figures irréductibles. \deleted{\unsure{Ceci}} (NB Dans les cas les plus importants, toute figure est finie, et a fortiori de type fini.)

Ceci dit, si $\mathcal F_{\mathrm{t.f.}}$ désigne l'ens figures de type fini, on récupère $\overline{\mathcal F_{\mathrm{t.f}}}$, l'ens des parties correspondantes de $M$, par la seule connaissance de $(M,\leq,R)$,
\scanpage{61}\mspage{60}
de la façon suivante : ce sont les $A\unsure{\in|\subset}M$ tels que $A=\bigcup_i A_{x_i}$, où \deleted{les} $\{x_i\}$ \deleted{sont}\inserted{est} un ens fini d'éléments de $M$, mutuellement compatibles (\deleted{\unsure{pour}}\inserted{au sens de} $R$.).

Donc une structure $(\mathcal F,\leq)$ satisfaisant $\deleted{\unsure{C_1}}C1,\deleted{\unsure{C_2}}C2,\deleted{\unsure{C_3}}C3$ + figures \deleted{\unsure{toutes|dont}} de type fini, équivaut à $(M,\leq,R)$, où \inserted{$R$ \uline{rel. sym. refl.}} \deleted{\unsure{$(M,\leq)$}}\inserted{où $\leq$ et $R$} satisfont \deleted{deux} conditions.

a) \inserted{(concernant $\leq$)} Pour tte partie fermée \uline{majorée} $A$\deleted{\unsure{,}} de $M$, $A$ est de type fini i.e de la forme $\bigcup_{\deleted{\unsure{x_i}}i}\deleted{\unsure{A_{x_i}}}M_{\leq\unsure{x_i|x}}$, où les $x_i\in A$ \deleted{\unsure{fo}} sont \unsure{en|un} \unsure{nb|ens} fini [ pour ceci, il suffit que les $M_{\leq x}$ soient finis ]

% slanted
\margin{donc toute partie fermée d'une partie fermée de type fini est de type fini}

b) (concernant $\leq$ \uline{et} $R$) Si $(x,y)\in R$, et $x'\leq x$ $y'\leq y$, alors $\deleted{\unsure{R}}(x',y')\in R$.

% slanted
\margin{C'est la \unsure{propriété|prop.} \unsure{Cond|Const} 2 de page 33}

(Les ``figures'' de $M$ sont alors les parties \deleted{fermées}\inserted{$A$} de $M$ qui ont les propriétés suivantes :
\[
\begin{aligned}
\text{a) } & A\text{ est fermé}\\
\text{b) } & A\text{ est de type fini}\\
\text{c) } & \deleted{\unsure{A}}\forall\,x,y\in A,\ \text{on a}\quad (x,y)\in R.
\end{aligned}
\]
\scanpage{62}\mspage{61}
\deleted{\unsure{Soit}} Ainsi, les \unsure{systèmes|système} $(\mathcal F,\leq)$ satisfaisant $C1,C2,C3$ et \unsure{tels|telle} que toute $F\in\mathcal F$ soit ``finie'' (i.e. l'ens $\mathcal F_{\leq F}$ des sous-figures est fini, ou encore, l'ens $M_{\leq F}\inserted{=\overline F}$ \deleted{des} de \deleted{\unsure{sous}} ses strates est fini) équivaut \unsure{:|à}

$(M,\leq,R)$ \deleted{\unsure{str.}} où $\leq$ relation d'ordre, $R$ relation refl. et sym., satisfaisant
\[
\begin{aligned}
\text{a) }&\text{les }M_{\leq x}\text{ finis}\quad(x\in M)\\
\text{b) }&\text{Si }(x,y)\in R\text{ et }x'\leq x,y'\leq y,\ \text{on a }\deleted{\unsure{R}}(x',y')\in R\,;
\end{aligned}
\]
où les figures sont interprétées comme les parties \uline{fermées}, \uline{finies} \inserted{$A$} de $M$, telles que $\forall x,y\in A$, on ait $(x,y)\in R$.

% slanted
\margin{\unsure{Ainsi} la notion \uline{combinatoire} de connexité. Elle \illegible{1} en évidence \deleted{que} aussi la notion \illegible{1} qui la sous-entend, que les strates $x$\unsure{,} sont 0-connexes + \ldots\ En fait, si on \unsure{voulait} \deleted{``exprimer''} \unsure{$=$|--} que les \unsure{conditions} \unsure{pr.} $A$ \unsure{clos|closes} \deleted{\unsure{et}} de $x$ \unsure{soient|soit} \unsure{co}}

\uline{Composantes connexes} Une figure est dite \uline{connexe} si \deleted{\unsure{on a}}
\[
F=F'\cup F''\quad
\inserted{\text{avec }F',F''\ \text{\uline{disjoints} i.e. }F'\cap F''=\varnothing}
\ \Longrightarrow\ F'=\deleted{\unsure{\varnothing}}F\ \text{ou }F''=\deleted{\unsure{\varnothing}}F.
\]
\margin{\uline{NB} $F$ \unsure{connexe} \unsure{ssi|ss} $\overline F$ l'est \unsure{comme ens.} \unsure{ordonné|préordonné}}

% slanted
\margin{On écrit $F=F'\amalg F''$ (signe de la somme disjointe\unsure{.|)} (\uline{NB} Si \deleted{\unsure{deux}} \unsure{parmi} les \deleted{\unsure{fig.}} $F',F'',F$ sont connexes, la troisième aussi.}

Une figure est dite \uline{0-connexe} si elle est connexe et $\ne\varnothing_{\mathcal F}$. Une sous-figure $F_0$ de $F$ est appelée \uline{composante connexe} de $F$ si a) elle est 0-connexe et

b) \deleted{$F_0$ est une \unsure{compos}} \inserted{$F=F_0\amalg G$, $G$ \unsure{compatible|composible|composable}} \deleted{\unsure{avec} $F$, disjoint de $F_0$.}

\deleted{\unsure{Il est}} \deleted{\uline{NB}}\inserted{Lemme} Soit $F=F'\amalg F''$, \deleted{avec $F'\cap F''=\varnothing$}, et soit $G$ une sous-figure connexe de $G$. Alors $G$ est contenue ds $F'$ ou $F''$.
\scanpage{63}\mspage{62}
En effet, $G=G'\amalg G''$ \unsure{ou|où} $G'\subset F',G''\subset F''$, \deleted{$G'$ et $G''$ disjoints,} \inserted{(cf p. 57)} donc on a $G'=G$ ou $G''=G$. donc $G\subset F'$ ou $G\subset F''$ q.e.d.

\uline{Corollaire} 1 \deleted{Deux} Soient $F_0,F_1$ deux comp. connexes de $F$. Alors on a $F_0=F_1$ ou $F_0$ et $F_1$ sont disjoints, et \deleted{\unsure{alors}} \inserted{dans ce cas} $F=F_0\amalg F_{\unsure{1|2}}\amalg \deleted{\unsure{G|F}}H$

On écrit $F=F_0\amalg G$, et on applique \unsure{lemme}

\deleted{\uline{Corollaire 2} Soit $(F_i)_{\unsure{j|i}\in I}$ une famille de composantes connexes deux à deux disjointes. Alors} \deleted{$\unsure{F=\coprod F_i}$} \deleted{la figure engendrée par les $F_i$ est leur somme disjointe $\coprod F_i$} \deleted{\unsure{et}} \deleted{$\subset F$, et si la famille est finie (p. ex. $F$} \deleted{\unsure{finie}} \deleted{figure finie) alors $F=(\coprod F_i)\amalg H$.}

\deleted{Il n'est pas évident que tte figure non vide ait une composante connexe,} \deleted{la difficulté étant dans l'\unsure{existence} \unsure{ou}} \deleted{\unsure{à} des sous-figures \unsure{``complém.res''|``comprises''|``compl. disj.''}} \deleted{\inserted{Mais c'est vrai pour les figures de t.f., et a fortiori pour les figures finies. Les}}

\uline{Proposition} Soit $F$ une figure, $A=\overline F\inserted{\subset M}$ l'ens \unsure{des|de} ses strates. \deleted{Alors l'application} Pour qu'une sous-figure \deleted{$\unsure{F|F_0|F_i}$} $F_0$ de $F$ soit une comp. conn. de $F$, il f et s que $\overline{F_0}$ soit une composante connexe de l'ens ordonné $A$ (pour l'ordre induit par $\leq$). On trouve ainsi une corr. biunivoque entre l'ens. des comp. conn.
\scanpage{64}\mspage{63}
de $F$, et celles de $A$. De plus, $F$ est somme directe disjointe de ses comp. connexes, i.e. $F=\mathop{\mathrm{Sup}}\limits_i F_i$, les comp. connexes $F_i$ deux à deux disjointes\ldots .

% slanted
\margin{\uline{NB} Dans tout ceci, on n'a pas besoin de $C_3$}

\uline{Composantes connexes \unsure{du|de} \deleted{$\unsure{\mathcal F|F}$}\unsure{Sle|site|Syst.} entier $\mathcal F$.}

Soit $M=\coprod_{i\in I}M_i$ la décomposition de $M$ en composantes connexes.

Pour toute figure $F$, posons
\[
F_i=\mathop{\mathrm{Sup}}\limits_{x\in M_i\cap M_{\leq F}}x
\]
% slanted
\margin{On \unsure{appelle} $F_i$ la \deleted{\unsure{$F_i$}} $i$-composante de $F$, ou la trace de $F$ sur $M_i$ (ou sur $\mathcal F_i$)}

Alors $F$ est la somme directe des $F_i$, \deleted{\unsure{et ?}} Si $F$ est de t.f. \inserted{(ou $I=\pi_0(M)$ fini!)} alors tous les $F_i$ sauf un nb fini \unsure{sont} $\varnothing$.

\deleted{\unsure{Inscrivons}} Soit $\mathcal F_i$ l'ens des $F\in\mathcal F$ telle que $\overline F\subset M_i$. Alors \deleted{\unsure{si}} $\mathcal F_i$ satisfait les conditions $C1,C2$\deleted{$\unsure{,C3}$} (et $C3$ aussi, si $\mathcal F$ y satisfait), et $M_i$ est l'ens de ses objets irréductibles.

\deleted{et l'application}
\diagram{\begin{tikzcd}[ampersand replacement=\&,column sep=large,row sep=small,
execute at end picture={\draw[thin,rounded corners=3pt] (oldfamily.south west) -- (oldfigure.west) -- (oldtail.east) .. controls +(right:5pt) and +(right:6pt) .. (oldtail.north east) -- (oldproduct.south west);}]
|[alias=oldfamily]| \mathcal F \arrow[r] \& |[alias=oldproduct]| \prod\mathcal F_i \\
|[alias=oldfigure]| \deleted{F} \arrow[r,mapsto] \& |[alias=oldtail]| \deleted{(F_i)_{i\in I}\ \text{\unsure{ou|où} }F_i\ \text{\unsure{est le}}}
\end{tikzcd}}
Si $F\unsure{\in|\subset}\mathcal F$, si $F_i$ sont ses composantes \unsure{suiv.} les $M_i$, alors $\overline{F_i}=\overline F\cap M_i$, $\overline F=\deleted{\unsure{\coprod|\bigcup}}\mathop{\unsure{\bigcup|\coprod}}\limits_i\overline{F_i}$ (réunion disj.)

Pour $i$ fixé, l'application $F\longrightarrow F_i$ commute aux \inserted{Sup \unsure{quelc.}, et aux Inf non vides}

L'application
\diagram{\begin{tikzcd}[ampersand replacement=\&,column sep=large]
\mathcal F \arrow[r] \& \prod\mathcal F_i \& F \arrow[r,mapsto] \& (F_i)
\end{tikzcd}}
est injective, et la relation d'ordre de $\mathcal F$ est induite par l'ordre produit
\scanpage{65}\mspage{64}
Inversement, donnons nous, \deleted{$\unsure{F_i\in|F_i}$} $\unsure{(}\forall i\in I)$ un $F_i\in\mathcal F_i$, quand peut-on dire qu'il existe $F\in\mathcal F$ telle que $\forall i$, $F_i$ soit la $i$ composante de $\deleted{\unsure{\mathcal F}}\unsure{F_0|F}$ ?

Déjà, il le faudrait pour \unsure{le|les} cas où tous les $F_i$ sauf un nb fini sont \deleted{$\varnothing$} vides. On va supposer donc

\uline{$C4$} Si $x,y\in M$ appartiennent à \unsure{des|deux} comp. connexes disjointes de $M$, alors $x$ et $y$ sont compatibles.

\deleted{Cor} \inserted{C'est \unsure{néc|néc.} et suffisant pour la \deleted{\unsure{déd|val}} \unsure{validité|déduction} des}

\uline{Corollaire 1} L'image de $\mathcal F\to\prod\mathcal F_i$ contient les $(F_i)_{i\in I}$ tels que $F_i=\varnothing$ \deleted{sauf} pour \deleted{\unsure{pres}}que tout $i$.

\inserted{\uline{Cor. 2}} \deleted{\unsure{donc}} $\mathcal F\to\prod\mathcal F_i$ est une bijection dans \deleted{les} le \deleted{deux} cas où \deleted{suivants} \inserted{$I$ est} \deleted{\unsure{f}} fini.

\deleted{a) $I$ est fini}

\deleted{b) Toute figure est de t.f. \unsure{--} alors}

\uline{Cor 3} Supposons que les figures \deleted{\unsure{de}} $\in\mathcal F$ soient de t.f. \deleted{Alors il} \inserted{(resp. finies)}. Alors il en est de $\unsure{=|\text{même}}$ \deleted{\unsure{de}} pour $\mathcal F_i$, de plus $\mathcal F\unsure{\text{---}|\longrightarrow}\prod\mathcal F_i$ a comme image l'ens des $(F_i)$ tels que $F_i$ soit $\varnothing$ pour pres que tout $i$.
